\subsection{Deutsch--Jozsa y Simon: Estructura y Esparsidad}
\label{subsec:patrones_dj_simon}

Varios algoritmos tempranos construyen superposiciones altamente estructuradas mediante oráculos e interferencia.

En Deutsch--Jozsa, la interferencia inducida por el oráculo y la transformación final de Hadamard concentra la amplitud sobre un subconjunto reducido de estados base cuando la función pertenece a una de las clases consideradas por el algoritmo. Este tipo de concentración implica esparsidad o repetición exacta de amplitudes, lo que favorece la compresión sin pérdida \parencite{DeutschJozsa1992, Nielsen2010}.

En Simon, tras medir el registro de salida, el registro de entrada colapsa a una superposición de dos estados relacionados por una máscara secreta. Ese estado intermedio es extremadamente esparso y, por tanto, muy compresible. Posteriormente, la aplicación de compuertas de Hadamard produce una superposición uniforme sobre un subconjunto definido por una restricción lineal, donde reaparecen ceros exactos y amplitudes repetidas \parencite{Simon1997, Nielsen2010}.
